\begin{abstract}
Machine unlearning for privacy-critical biometric systems, fingerprint, iris, and face recognition, remains underexplored despite GDPR/CCPA mandates requiring erasure throughout deployment. Existing methods target classification and generative tasks, collapsing all forget classes to a single distant point, which inflates convergence time on resource-constrained edge devices facing frequent, minimal-replay removal requests; they are also benchmarked almost exclusively on ResNet/ViT for CIFAR/ImageNet, leaving biometric backbones untested. We propose Camouflage Unlearning (CU), which maps forget classes to nearby retain-class boundaries instead of a single point, reducing migration distance; this proximity target shapes only the training-time optimization objective, and we verify empirically that it does not induce elevated impostor matches against the assigned retain identity at verification time. CU achieves 3--4$\times$ faster convergence via retain-dominance constraints and localized regularization that surgically updates specific transformer blocks, while maintaining competitive retain accuracy. Across classification, iris, fingerprint, and face recognition, CU shows competitive performance and privacy behavior (MIA $\approx$ 0.5), supporting its effectiveness for regulatory-compliant, deep embedding-based biometric authentication.
\end{abstract}

\begin{IEEEkeywords}
Biometric recognition, data privacy, face recognition, GDPR, machine unlearning, vision transformers
\end{IEEEkeywords}